\documentclass[11pt]{article}
\usepackage{amsmath}

\setlength{\parindent}{0pt}
\setlength{\parskip}{0.7em}
\setlength{\textwidth}{6.5in}
\setlength{\oddsidemargin}{0in}
\setlength{\evensidemargin}{0in}
\setlength{\topmargin}{-0.4in}
\setlength{\textheight}{9.2in}

\begin{document}

\begin{center}
{\Large SYSC 5108 Exam Study: Assignment 2, Question 5}\\
\vspace{0.3em}
{\large Logistic regression gradient for $0/1$ and $-1/+1$ label encodings}\\
\vspace{0.3em}
Source question: \texttt{SYSC 5108 W26 Assignment 2.pdf}, Question 5
\end{center}

\section*{Question}

Consider training a logistic regression model with $K=2$ classes. First assume $y\in\{0,1\}$ and
\[
J(w)=-y\log(\sigma(z))-(1-y)\log(1-\sigma(z)),
\qquad
z=w^T x,\qquad x_0=1.
\]
Show that
\[
\nabla_w J(w)=-x\bigl(y-\sigma(z)\bigr).
\]
Then change the encoding to $y\in\{-1,1\}$, show the result again, and explain the impact.

\section*{Step 1: Recall the sigmoid derivative}

Chapter 3 slides 32--34 define the logistic sigmoid:
\[
\sigma(z)=\frac{1}{1+e^{-z}}.
\]
The derivative that we need is
\[
\frac{d\sigma(z)}{dz}=\sigma(z)(1-\sigma(z)).
\]
Also, since
\[
z=w^T x,
\]
the gradient of $z$ with respect to $w$ is
\[
\nabla_w z=x.
\]

\section*{Part (a): Derivation for $y\in\{0,1\}$}

Let
\[
s=\sigma(z).
\]
The loss is
\[
J(w)=-y\log(s)-(1-y)\log(1-s).
\]
First differentiate with respect to $s$:
\[
\frac{dJ}{ds}
=-\frac{y}{s}+\frac{1-y}{1-s}.
\]
Now use
\[
\frac{ds}{dz}=s(1-s).
\]
By the chain rule,
\[
\frac{dJ}{dz}
=
\frac{dJ}{ds}\frac{ds}{dz}
=
\left(-\frac{y}{s}+\frac{1-y}{1-s}\right)s(1-s).
\]
Simplify:
\[
\frac{dJ}{dz}
=
-y(1-s)+(1-y)s.
\]
Expand the terms:
\[
\frac{dJ}{dz}
=-y+ys+s-ys=s-y.
\]
Since $s=\sigma(z)$,
\[
\frac{dJ}{dz}=\sigma(z)-y.
\]
Finally apply the chain rule through $z=w^T x$:
\[
\nabla_w J(w)
=
\frac{dJ}{dz}\nabla_w z
=
(\sigma(z)-y)x.
\]
This is the same as
\[
\boxed{
\nabla_w J(w)
=
-x\bigl(y-\sigma(z)\bigr)
}.
\]

\section*{Step 2: Interpret the $0/1$ Gradient}

For $y=1$:
\[
\nabla_w J(w)=-(1-\sigma(z))x.
\]
Gradient descent uses
\[
w_{\text{new}}=w-\epsilon\nabla_wJ(w),
\]
so the update adds a positive multiple of $x$ when the positive class is underpredicted.

For $y=0$:
\[
\nabla_w J(w)=\sigma(z)x.
\]
Gradient descent subtracts a positive multiple of $x$, pushing the logit $z=w^Tx$ downward.

This matches the intuition: if the target is $1$, push the probability upward; if the target is $0$, push the probability downward.

\section*{Part (b): Change the Encoding to $y\in\{-1,1\}$}

When labels are encoded as $-1$ and $+1$, do not plug $y=-1$ directly into the $0/1$ cross-entropy formula. That formula assumes $y$ is an indicator variable. Instead, use the equivalent binary logistic loss written in margin form:
\[
J(w)=\log(1+e^{-yz}),
\qquad
y\in\{-1,1\}.
\]
Here,
\[
yz
\]
is the signed margin. If $yz$ is large and positive, the example is classified correctly with high confidence. If $yz$ is negative, the example is on the wrong side of the decision boundary.

Differentiate with respect to $z$:
\[
\frac{dJ}{dz}
=
\frac{1}{1+e^{-yz}}\cdot e^{-yz}(-y).
\]
So
\[
\frac{dJ}{dz}
=
-\frac{y e^{-yz}}{1+e^{-yz}}.
\]
Divide numerator and denominator by $e^{-yz}$:
\[
\frac{dJ}{dz}
=
-\frac{y}{1+e^{yz}}.
\]
Using
\[
\sigma(-yz)=\frac{1}{1+e^{yz}},
\]
we get
\[
\frac{dJ}{dz}=-y\sigma(-yz).
\]
Now use $\nabla_wz=x$:
\[
\boxed{
\nabla_w J(w)
=
-y\,x\,\sigma(-yz)
}.
\]
Equivalently,
\[
\boxed{
\nabla_w J(w)
=
-\frac{yx}{1+e^{yz}}
}.
\]

\section*{Step 3: Show Consistency with the $0/1$ Formula}

Let the old label be $t\in\{0,1\}$ and the new label be
\[
y=2t-1.
\]
If $t=1$, then $y=1$. The $-1/+1$ gradient becomes
\[
-y x\sigma(-yz)=-x\sigma(-z).
\]
Since
\[
\sigma(-z)=1-\sigma(z),
\]
this is
\[
-x(1-\sigma(z)),
\]
which matches the $0/1$ result for $t=1$:
\[
x(\sigma(z)-1).
\]

If $t=0$, then $y=-1$. The $-1/+1$ gradient becomes
\[
-(-1)x\sigma(z)=x\sigma(z),
\]
which matches the $0/1$ result for $t=0$:
\[
x(\sigma(z)-0)=x\sigma(z).
\]

So the two encodings describe the same binary logistic regression objective if the loss is written correctly for each encoding.

\section*{Step 4: Explain the Impact of Changing the Encoding}

Changing from $0/1$ labels to $-1/+1$ labels does not change the decision boundary. In both cases, the boundary is
\[
z=w^Tx=0.
\]
For the $0/1$ encoding,
\[
p(y=1\mid x)=\sigma(z),
\]
and the usual threshold is $p\ge 0.5$, which is equivalent to $z\ge 0$.

For the $-1/+1$ encoding, the loss is written using the margin $yz$. This form is often cleaner because:
\begin{itemize}
\item Correct and confident examples have large positive $yz$ and therefore a small loss.
\item Incorrect examples have negative $yz$ and therefore a large loss.
\item The gradient magnitude is controlled by $\sigma(-yz)$, so examples that are badly misclassified receive a larger update.
\end{itemize}

The main caution is that the old $0/1$ cross-entropy formula cannot be reused unchanged with $y=-1$. You must rewrite it as
\[
J(w)=\log(1+e^{-yz}).
\]

\section*{Concise Exam Answer}

For $y\in\{0,1\}$, let $s=\sigma(z)$. Then
\[
J(w)=-y\log(s)-(1-y)\log(1-s).
\]
The chain rule gives
\[
\frac{dJ}{dz}
=
\left(-\frac{y}{s}+\frac{1-y}{1-s}\right)s(1-s)
=s-y
=\sigma(z)-y.
\]
Since $\nabla_w z=x$,
\[
\nabla_w J(w)=(\sigma(z)-y)x=-x(y-\sigma(z)).
\]
For $y\in\{-1,1\}$, use
\[
J(w)=\log(1+e^{-yz}).
\]
Then
\[
\frac{dJ}{dz}=-\frac{y}{1+e^{yz}}=-y\sigma(-yz),
\]
so
\[
\nabla_wJ(w)=-yx\sigma(-yz)=-\frac{yx}{1+e^{yz}}.
\]
The impact is mostly notational and algebraic: $-1/+1$ labels express the loss in terms of the signed margin $yz$, while the $0/1$ form expresses it as Bernoulli cross-entropy. The decision boundary remains $w^Tx=0$ when the loss is written consistently.

\section*{Cross References Used}

\begin{itemize}
\item Assignment: \texttt{SYSC 5108 W26 Assignment 2.pdf}, Question 5.
\item Local submitted Assignment 2 PDF checked for comparison, Question 5 response.
\item Slides: Chapter 3 slide deck, slides 32--34 on the logistic sigmoid and binary logistic regression.
\item Slides: Chapter 3 slide deck, logistic regression training slide on log-likelihood and negative log-likelihood.
\item Textbook: Goodfellow, Bengio, and Courville, \textit{Deep Learning}, Chapter 3, Section 3.10 on the logistic sigmoid.
\item Textbook: Goodfellow, Bengio, and Courville, \textit{Deep Learning}, Chapter 6, Section 6.2.2.2 on sigmoid output units for Bernoulli output distributions and maximum likelihood.
\end{itemize}

\end{document}
